\documentclass[10pt,a4paper]{amsart}
\usepackage[margin=19mm]{geometry}
\usepackage{amsmath,amssymb,mathtools}
\usepackage[scr=rsfs,cal=euler]{mathalfa}
\usepackage{xfrac}
\usepackage[unicode,hidelinks,bookmarks=false]{hyperref}
\newcommand{\ie}{i.e.\ }
\newcommand{\cf}{cf.\ }
\newcommand{\resp}{respectively\ }
\newcommand{\Subsection}{\subsection}
\providecommand{\nobreakdash}{\nobreak\hbox{-}}
\setlength{\emergencystretch}{2em}
\multlinegap0pt
\usepackage{tikz}
\usetikzlibrary{cd}
\usepackage{amssymb}
\usepackage{enumerate}
\usepackage[all]{xy}
\usepackage{float}
\usepackage{mathtools}
\usepackage{caption}
\usepackage{xcolor}
\usepackage[normalem]{ulem}
\usepackage[shortlabels]{enumitem}
\usepackage{booktabs}
\newtheorem{assumption}{Assumption}
\newtheorem{lemma}{Lemma}
\newcommand{\EC}{\mathsf{E}}
\newcommand{\Q}{\mathbb{Q}}
\newcommand{\Sel}{\operatorname{Sel}}
\title{Iwasawa Theory of Elliptic Curves in Quadratic Twist Families}
\author{Debanjana Kundu, Katharina M\"uller}
\date{Source: arXiv:2507.21339v4 (CC BY 4.0)}
\begin{document}
\maketitle

\noindent\emph{Source and licence.} Iwasawa Theory of Elliptic Curves in Quadratic Twist Families. Version arXiv:2507.21339v4 (CC BY 4.0), distributed by arXiv under the Creative Commons Attribution 4.0 International licence, http://creativecommons.org/licenses/by/4.0/, as recorded in the arXiv licence metadata for this exact version and shown on the abstract page https://arxiv.org/abs/2507.21339v4. Full licence notice: \texttt{licenses/arxiv-2507.21339-CC-BY-4.0.txt}.\par\smallskip

\noindent\textit{Editorial scope.} Counted unit: the article's lemma isom:different-selmer (source0.tex lines 837--845). The article's complete original proof of it is delivered with it (source0.tex lines 847--859, one \texttt{proof} environment of 13 lines). No mathematics is written, repaired or re-derived: the statement and the proof are the article's own lines, in the article's own notation, and no retained line is substituted or re-flowed.  Retained uncounted as the source-local context the proof needs: lines 751--760.  Nothing was omitted from any retained span, so the package carries no omission record.  No retained line contains a citation, so the wrapper carries no bibliography and no bibliography entry.  The retained lines contain the article's own diagram code, which the wrapper typesets with the standard tikz packages; no image file is delivered and the package declares no asset dependency.  Editorial formatting only, in addition: an AMS \texttt{amsart} preamble that restates the article's own declarations and macros used by the retained lines, namely the environments \texttt{assumption}, \texttt{lemma} on separate counters, together with the standard packages those lines need.  The retained environments print in this document as Assumption 1, Lemma 1 in order of appearance, whereas the article prints the same items with the numbering of its own full document.  The complete native source of the article as distributed in the arXiv e-print for this exact version is bundled unchanged as \texttt{sources/182340/source0.tex}.
\begin{assumption}\label{ass}Let $K$ be a quadratic field such that the following conditions are satisfied
    \begin{enumerate}
    \item[\textup{(}i\textup{)}] $\EC(K_v)[p]=\{0\}$ for all $v\in \Sigma(K)$.
    \item[\textup{(}ii\textup{)}] $\EC$ has good ordinary reduction at $p$ odd.
    \item[\textup{(}iii\textup{)}] $p$ splits in $K$.
    \item [\textup{(}iv\textup{)}] $\Sel^0(\EC/K)= \{0\}$.
    \item[\textup{(}v\textup{)}]
    Assume that $\Sel(\EC/\Q)$ and $\Sel(\EC^K/\Q)$ are both non-trivial.
\end{enumerate}
\end{assumption}

\begin{lemma}
\label{isom:different-selmer}
Suppose that Assumption~\ref{ass}\textup{(}i\textup{)} is satisfied\footnote{we do not insist that the conditions (ii)--{(v)} hold.}.
Let $k$ be a positive integer.
Then
\[
\Sel^0(\EC/K)[p^k]\cong \Sel^0(\EC[p^k]/K).
\]
\end{lemma}

\begin{proof}
Consider the following commutative diagram:
\[
\begin{tikzcd}
        0\arrow[r] &\Sel^0(\EC[p^k]/K) \arrow[r]\arrow[d]&H^1(G_\Sigma(K),\EC[p^k])\arrow[r]\arrow[d,"h"]& \prod_{v\in \Sigma(K)}H^1(K_{v},\EC[p^k])\arrow[d,"g"]\\
        0\arrow[r] &\Sel^0(\EC/K)[p^k] \arrow[r]&H^1(G_\Sigma(K),\EC[p^\infty])[p^k]\arrow[r]& \prod_{v\in \Sigma(K)}H^1(K_{v},\EC[p^\infty])[p^k]
\end{tikzcd}
\]
The maps $h,g$ are the usual maps that arise in the definition of the Kummer sequence.
Recall that $\ker(h) = H^0(G_\Sigma(K),\EC[p^\infty])/p^k$ and $\ker(g)=\prod_{v\in \Sigma(K)}H^0(K_v,\EC[p^\infty])/p^k$ but both of the cohomology groups are trivial by Assumption~\ref{ass}(i).
Furthermore, $h$ is surjective by definition.
The claim now follows from the snake lemma.
\end{proof}

\end{document}
